% Part:sets-functions-relations
% Chapter: size-of-sets
% Section: reduction-alt

\documentclass[../../../include/open-logic-section]{subfiles}

\begin{document}

\olfileid{sfr}{siz}{red-alt}

\olsection{归约}

\begin{editorial}
  本节通过归约证明不可数性，其结果与 \olref[nen-alt]{sec} 相符。另一种与 \olref[nen]{sec} 结果相匹配、稍加详细的版本在 \olref[red]{sec} 给出。
\end{editorial}

我们通过对角化论证证明了 $\Bin^\omega$ 不可数。我们也用类似的对角化论证证明了 $\Pow{\Nat}$ 不可数。但证明 $\Pow{\Nat}$ 不可数还有另一种方法：证明\emph{如果 $\Pow{\Nat}$ 可数，则 $\Bin^\omega$ 也可数}。既然我们知道 $\Bin^\omega$ 不可数，那么 $\Pow{\Nat}$ 也就不可数。

这种做法称为将一个问题\emph{归约}到另一个问题。在本例中，我们将枚举 $\Bin^\omega$ 的问题归约为枚举 $\Pow{\Nat}$ 的问题。后者的一个解（即 $\Pow{\Nat}$ 的一个枚举）将给出前者的一个解（即 $\Bin^\omega$ 的一个枚举）。

要将枚举集合~$B$ 的问题归约为枚举集合~$A$ 的问题，我们需要提供一种方法，将~$A$ 的枚举转化为~$B$ 的枚举。实现这一目标最简单的方法是定义一个满射 $f\colon A \to B$。如果 $x_1$, $x_2$, \dots{} 枚举了~$A$，那么 $f(x_1)$, $f(x_2)$, \dots{} 就会枚举~$B$。在我们的例子中，我们要寻找一个满射 $f\colon \Pow{\Nat} \to \Bin^\omega$。

\begin{prob}
证明：如果存在单射函数 $g\colon B \to A$，且 $B$~是不可数的，那么 $A$~也是不可数的。请通过说明如何利用~$g$ 将~$A$ 的枚举转化为~$B$ 的枚举来完成证明。
\end{prob}

\begin{proof}[通过归约法证明 {\olref[nen-alt]{thm:nonenum-pownat}}]
为进行归约，假设 $\Pow{\Nat}$ 是可数的，从而存在它的一个枚举 $N_{1}$, $N_{2}$, $N_{3}$, \dots。

定义函数 $f \colon \Pow{\Nat} \to \Bin^\omega$，令 $f(N)$ 为满足以下条件的字符串 $s_{k}$：$s_{k}(n) = 1$ 当且仅当 $n \in N$，否则 $s_k(n) = 0$。

这显然定义了一个函数，因为对于任意 $N \subseteq \Nat$，任何 $n \in \Nat$ 要么属于 $N$，要么不属于 $N$。例如，偶自然数集合 $2\Nat = \Setabs{2n}{n \in \Nat} = \{0,2, 4, 6,
\dots\}$ 映射到字符串 $1010101\dots$；$\emptyset$ 映射到 $0000\dots$；$\Nat$ 映射到 $1111\dots$。

它也是满射：每一个由 $0$ 和 $1$ 组成的字符串都对应于某个自然数集合，即由字符串中包含 $1$ 的位置所对应的自然数组成的集合。更准确地说，如果 $s \in \Bin^\omega$，则如下定义 $N \subseteq \Nat$：
\[
N = \Setabs{n \in \Nat}{s(n) = 1}
\]
那么 $f(N) = s$，对照 $f$ 的定义即可验证。

现在考虑列表
\[
f(N_1), f(N_2), f(N_3), \dots
\]
由于 $f$ 是满射，$\Bin^\omega$ 中的每个成员必定作为 $f$ 在某个自变元处的值出现，从而必定出现在该列表中。因此这个列表必然枚举了所有的 $\Bin^\omega$。

所以，如果 $\Pow{\Nat}$ 是可数的，那么 $\Bin^\omega$ 将是可数的。但 $\Bin^\omega$ 是不可数的（\olref[nen-alt]{thm:nonenum-bin-omega}）。因此 $\Pow{\Nat}$ 是不可数的。
\end{proof}

%\begin{explain}
%It is easy to be confused about the direction the reduction goes in.
%For instance, !!a{surjective} function $g \colon \Bin^\omega \to X$
%does \emph{not} establish that $X$ is !!{nonenumerable}.  (Consider $g
%\colon \Bin^\omega \to \Bin$ defined by $g(s) = s(1)$, the function
%that maps a sequence of $0$'s and $1$'s to its first !!{element}.  It
%is surjective, because some sequences start with $0$ and some start
%with $1$. But $\Bin$ is finite.)  Note also that the function $f$ must
%be surjective, or otherwise the argument does not go through:
%$f(x_1)$, $f(x_2)$, \dots{} would then not be guaranteed to include
%all the !!{element}s of~$Y$. For instance, $h\colon \Nat \to
%\Bin^\omega$ defined by
%\[
%h(n) = \underbrace{000\dots0}_{\text{$n$ $0$'s}}
%\]
%is a function, but $\Nat$ is !!{enumerable}.
%\end{explain}

\begin{prob}\label{sfr:siz:red:prob:nat-nat}
  用归约论证证明所有函数 $f\colon \Nat \to \Nat$ 的集合 $X$ 是不可数的。（提示：给出一个从 $X$ 到 $\Bin^\omega$ 的满射函数。）
\end{prob}

\begin{prob}
用归约论证证明所有由自然数对构成的集合的集合，即 $\Pow{\Nat \times \Nat}$，是不可数的。
\end{prob}

\begin{prob}
用归约论证证明 $\Nat^\omega$，即自然数的无限序列的集合，是不可数的。
\end{prob}

%\begin{prob}
%Let $P$ be the set of functions from $\Nat$ to the set $\{0\}$, and let $Q$ be the set of \emph{partial}
%functions from the set of positive integers to the set $\{0\}$. Show
%that $P$~is !!{enumerable} and $Q$~is not. (Hint: reduce the problem
%of enumerating $\Bin^\omega$ to enumerating~$Q$).
%\end{prob}

\begin{prob}
设 $S$ 为所有从 $\Nat$ 到集合 $\{0,1\}$ 的满射的集合，即 $S$ 由所有满射 $f \colon \Nat \to \Bin$ 组成。证明 $S$ 是不可数的。
\end{prob}

\begin{prob}
证明所有实数构成的集合 $\Real$ 是不可数的。
\end{prob}

\end{document}